\begin{figure*}[t]
\centering
\begin{tikzpicture}[
    transform shape, % Scales nodes along with coordinates
    scale=0.9,       % Adjust this value (e.g., 0.85 to 0.95) to dial in the perfect fit
    node distance=0.4cm and 0.5cm, % Slightly tightens the global spacing between nodes
    font=\small,
    nbox/.style={draw, rounded corners=2pt, minimum height=0.8cm, align=center, inner sep=4pt},
    frozen/.style={nbox, fill=gray!12, draw=gray!60},
    learn/.style={nbox, fill=blue!10, draw=blue!60},
    nout/.style={nbox, fill=green!12, draw=green!55!black},
    cache/.style={nbox, fill=yellow!18, draw=orange!70},
    arr/.style={-{Latex[length=2mm]}, thick}
]
% Inputs
\node[frozen] (xt) at (0, 0) {Raw features\\$\mathbf{x}_t \in \mathbb{R}^p$};
\node[frozen, below=0.5cm of xt] (ct) {Context vector\\$\mathbf{s}_t \in \mathbb{R}^{d_c}$};

% Tokenization
\node[learn, right=0.9cm of xt] (tok) {Feature\\tokenizer\\$\mathbf{t}_j = a_j x_{t,j} + \mathbf{e}_j$};

% Context token
\node[learn, right=0.9cm of ct] (ctok) {Context\\token\\$\mathbf{c}_t = W_c \mathbf{s}_t$};

% Transformer block
\node[learn, right=1.0cm of tok, minimum width=3.6cm, minimum height=2.2cm] (tr) {Feature-Axis\\Transformer\\$L$ layers, $h$ heads};

% Pool
\node[learn, right=0.6cm of tr] (pool) {CLS\\readout\\$\mathbf{u}_t$};

% Constraint head
\node[learn, right=0.6cm of pool, minimum width=2.6cm] (head) {Combination Head\\(uncon./affine/convex)};

% Output weights
\node[nout, right=0.6cm of head] (wt) {$\mathbf{w}_t \in \mathbb{R}^K$};

% Base predictors
\node[cache, below=1.5cm of head] (yhat) {Cached base predictions\\$\hat{\mathbf{y}}_t = (\hat{y}^{(1)}_t,\dots,\hat{y}^{(K)}_t)$};

% Final
\node[nout, right=0.6cm of wt] (final) {$\tilde{y}_t = \mathbf{w}_t^{\!\top}\hat{\mathbf{y}}_t$};

% Loss
\node[nbox, fill=red!10, draw=red!60, below=0.5cm of final] (loss) {$\ell(y_t, \tilde{y}_t)$};

% Arrows
\draw[arr] (xt) -- (tok);
\draw[arr] (ct) -- (ctok);
\draw[arr] (tok) -- (tr.west |- tok);
\draw[arr] (ctok) -- (tr.west |- ctok);
\draw[arr] (tr) -- (pool);
\draw[arr] (pool) -- (head);
\draw[arr] (head) -- (wt);
\draw[arr] (wt) -- (final);
\draw[arr] (yhat) -| (final);
\draw[arr] (final) -- (loss);

% Backprop dashed arrow
% \draw[arr, dashed, red!70!black, bend right=20] (loss.west) to node[midway, below, font=\scriptsize]{$\nabla_\theta$ end-to-end} (tr.south);
\draw[arr, dashed, red!70!black] (loss.south) -- ++(0,+0.2cm) -| node[pos=0.25, below, font=\scriptsize]{$\nabla_\theta$ end-to-end} (tr.south);

\end{tikzpicture}
\caption{Overview of the proposed TAFS framework. Each of the $p$ engineered features is treated as a token with a learnable identity embedding, and a context token derived from lagged targets and regime indicators is prepended. A feature-axis transformer produces a task-aligned summary $\mathbf{u}_t$, which is mapped by a combination head to weights $\mathbf{w}_t$ over $K$ pretrained base forecasters under one of three constraint regimes (unconstrained, affine, convex). The final prediction is a weighted sum of cached base outputs $\hat{\mathbf{y}}_t$, and the cumulative prediction error is back-propagated end-to-end through the transformer and the head; base predictors are kept frozen.}
\label{fig:architecture}
\end{figure*}
